\documentclass[11pt,a4paper]{article}

% ── Packages ──
\usepackage[utf8]{inputenc}
\usepackage[T1]{fontenc}
\usepackage{amsmath,amssymb}
\usepackage{graphicx}
\usepackage[margin=1in]{geometry}
\usepackage{natbib}
\usepackage{hyperref}
\usepackage{booktabs}
\usepackage{xcolor}
\usepackage{siunitx}
\usepackage{caption}
\usepackage{lineno}

\linenumbers

% ── Macros ──
\newcommand{\vmr}{\mathrm{VMR}}

% ── Title ──
\title{Biospatial Regime Space: VMR Fingerprint Descriptors Resolve\\a Clustering Manifold Across Spatial Organization Mechanisms}

\author{
  Pretham Voruganti\textsuperscript{1}
  \\[0.5em]
  \small\textsuperscript{1}Independent Researcher
}

\date{}

\begin{document}

\maketitle

% ══════════════════════════════════════════════════════════════
\begin{abstract}
Multiscale variance-to-mean ratio (VMR) fingerprints detect spatial
organization across biological scales, but the question of which
\emph{generative mechanisms} produce distinguishable fingerprints---and
which do not---has remained open. Here we address this by simulating
seven canonical spatial organization mechanisms (random, clustered,
exclusion, layered, compartmental, branch-constrained, hierarchical)
with continuous parameter sweeps (3,750 total simulated patterns, 50
realizations per parameter setting), extracting six VMR-derived
fingerprint descriptors, and constructing a PCA-based regime space.

We find that VMR descriptors resolve a structured \emph{clustering
manifold}: branch-constrained, hierarchical, compartmental, and
clustered regimes occupy distinct, separable regions of descriptor space
(three PCs explain 96.0\% of variance). In contrast, exclusion, layered,
and random regimes collapse into a shared \emph{null cloud}, indicating
that VMR fingerprints---while powerful for detecting clustering-positive
mechanisms---do not distinguish sub-Poisson spatial organization. This
is an honest boundary of the method, not a failure: the clustering
manifold is real and stable.

Projecting 72 empirical fingerprints from three biological systems---protein
nanoclusters (SMLM, nm scale), synaptic organization (dendrites,
\si{\micro\meter} scale), and cell type architecture (tissue,
\si{\milli\meter} scale)---into this simulation-defined space, we find
that empirical data land in biologically interpretable regions: SMLM
near clustered, dendrites near layered, and tissue distributed across
multiple regimes. 44/72 empirical fingerprints are flagged as
out-of-regime, suggesting that real biological organization often
exceeds the complexity of single-mechanism simulations.
\end{abstract}

\clearpage

% ══════════════════════════════════════════════════════════════
\section{Introduction}
\label{sec:intro}

Spatial point pattern analysis has a long history in ecology, epidemiology,
and materials science \citep{Diggle2003, Baddeley2015, Ripley1977}. Recently,
multiscale variance-to-mean ratio (VMR) fingerprints have been applied to
biological systems at scales ranging from protein nanoclusters (nm) to cell
type organization in tissue (mm), providing a unified framework for detecting
and characterizing spatial structure.

A natural question arises: \emph{which spatial organization mechanisms produce
distinguishable VMR fingerprints, and which are degenerate?} If two
fundamentally different mechanisms produce indistinguishable fingerprints, the
method cannot resolve them---and users should know this.

Here we take a simulation-first approach. We define seven canonical spatial
organization mechanisms spanning the space of common biological patterns,
simulate each with continuous parameter sweeps to create regime
\emph{clouds} rather than single points, extract fingerprint descriptors,
and ask: does a stable, interpretable regime space emerge?

The answer is nuanced. VMR descriptors resolve a structured manifold of
clustering-positive mechanisms (clustered, compartmental, hierarchical,
branch-constrained) but place sub-Poisson regimes (exclusion, layered,
random) in a shared null cloud. We present this as an honest characterization
of what VMR fingerprints can and cannot distinguish.


% ══════════════════════════════════════════════════════════════
\section{Methods}
\label{sec:methods}

\subsection{Simulated Spatial Mechanisms}

All simulations operate on a $10{,}000 \times 10{,}000$ arbitrary unit ROI
with 2,000 points. Seven mechanisms are simulated with parameter sweeps:

\begin{enumerate}
  \item \textbf{Random}: Uniform Poisson process. Sweep: point density
    (500--10,000 points).
  \item \textbf{Clustered}: Thomas process \citep{Thomas1949}. Sweep:
    cluster radius (50--2,000 a.u.).
  \item \textbf{Exclusion}: Simple sequential inhibition (SSI) via KD-tree.
    Sweep: minimum inter-point distance (10--400 a.u.).
  \item \textbf{Layered}: Gaussian bands along one axis. Sweep: layer
    width fraction (0.1--1.0).
  \item \textbf{Compartmental}: SSI-placed compartment centers with
    Thomas-process sub-clustering. Sweep: number of compartments (2--25).
  \item \textbf{Branch-constrained}: L-system branching tree with Gaussian
    jitter. Sweep: branch width (5--500 a.u.).
  \item \textbf{Hierarchical}: Three-level nested Thomas process. Sweep:
    inter-level scale ratio (2--20).
\end{enumerate}

Each parameter setting is simulated 50 times, yielding 3,750 total point
patterns.

\subsection{VMR Fingerprinting}

For each simulated pattern, we compute multiscale VMR fingerprints using a
$256 \times 256$ grid at 12 logarithmically spaced scales. A complete
spatial randomness (CSR) null of 20 Poisson realizations provides the
baseline. The VMR ratio curve is $\vmr_\text{ratio}(s) =
\vmr_\text{observed}(s) / \vmr_\text{null}(s)$.

\subsection{Fingerprint Descriptors}

Six primary descriptors are extracted from each $\log_2(\vmr_\text{ratio})$
curve:

\begin{enumerate}
  \item \textbf{Peak height}: $\max_s \log_2(\vmr_\text{ratio}(s))$
  \item \textbf{Threshold scale}: first normalized scale where
    $|\log_2(\vmr_\text{ratio})| > 0.5$
  \item \textbf{Width}: FWHM in log-scale decades
  \item \textbf{Small-scale slope}: slope over the first third of scales
  \item \textbf{Large-scale slope}: slope over the last third of scales
  \item \textbf{Integrated excess}: $\int \log_2(\vmr_\text{ratio})\,
    d(\log s)$
\end{enumerate}

Two supplementary descriptors---anisotropy (eigenvalue ratio of 2D
covariance) and multimodality (second peak prominence ratio)---are
computed but not used in the embedding.

\subsection{Regime Space Construction}

PCA is fitted on the 3,750 simulated descriptor vectors (standardized to
zero mean, unit variance). The first three principal components define the
regime space. Empirical fingerprints are projected into this space via the
same standardization and PCA transformation---they do not influence the
embedding.

Regime separability is quantified by silhouette score across all seven
regimes. Soft regime affinity is computed via Gaussian KDE fitted per
regime in the 3D PCA space, with affinities normalized to sum to 1.
Out-of-regime status is flagged when an empirical point's nearest-neighbor
distance to any simulated point exceeds the 95th percentile of
within-simulation NN distances.

\subsection{Empirical Data}

72 empirical VMR fingerprints are loaded from three biological systems:

\begin{itemize}
  \item \textbf{Tissue} (44 fingerprints): Allen Brain Cell Atlas MERFISH,
    cell type VMR ratios from isocortex, hippocampus, and thalamus.
  \item \textbf{SMLM} (4 fingerprints): single-molecule localization
    microscopy of synaptic receptors, nuclear pores, membrane domains,
    and a negative control.
  \item \textbf{Dendrites} (24 fingerprints): excitatory and inhibitory
    synaptic inputs across 12 cortical neuron types from MICrONS
    connectomics data.
\end{itemize}


% ══════════════════════════════════════════════════════════════
\section{Results}
\label{sec:results}

\subsection{The clustering manifold is real; the null cloud is degenerate}

PCA on 3,750 simulated descriptor vectors yields a clear structure
(Fig.~\ref{fig:phase_diagram}). PC1 (72.3\% variance) separates
clustering-positive regimes from the null cloud, loading primarily on
peak height, integrated excess, and threshold scale
(Fig.~\ref{fig:pca_loadings}). PC2 (16.0\%) captures scale structure
differences among clustering regimes. PC3 (7.7\%) adds minor additional
separation. Together, three components explain 96.0\% of total variance.

The overall silhouette score across all seven regimes is 0.172, reflecting
the degeneracy of the null cloud. Within the clustering manifold,
branch-constrained, hierarchical, compartmental, and clustered regimes
are well-separated (Fig.~\ref{fig:descriptor_distributions}). Random,
exclusion, and layered regimes overlap substantially.

\begin{figure}[htbp]
  \centering
  \includegraphics[width=\textwidth]{../figures/fig_phase_diagram.pdf}
  \caption{Biospatial regime space (PC1 vs PC2). Clustering-positive
  regimes (clustered, compartmental, hierarchical, branch-constrained)
  form a structured manifold, while sub-Poisson regimes (random, exclusion,
  layered) collapse into a null cloud (dashed ellipse). Empirical
  fingerprints from tissue (stars), SMLM (triangles), and dendrites
  (diamonds) are projected into the simulation-defined space.}
  \label{fig:phase_diagram}
\end{figure}

\begin{figure}[htbp]
  \centering
  \includegraphics[width=\textwidth]{../figures/fig_simulated_patterns.pdf}
  \caption{Example point patterns from each of the seven simulated regimes.}
  \label{fig:simulated_patterns}
\end{figure}

\subsection{Empirical projections are biologically interpretable}

Projecting 72 empirical fingerprints into the simulation-defined regime
space (Fig.~\ref{fig:phase_diagram}), we observe biologically coherent
placement:

\begin{itemize}
  \item \textbf{SMLM}: synaptic receptors, nuclear pores, and membrane
    domains project near the clustered regime; the negative control projects
    near random.
  \item \textbf{Dendrites}: most synaptic input fingerprints project near
    layered, consistent with distance-dependent synapse placement along
    dendritic branches. Some inhibitory inputs on specific cell types project
    near hierarchical.
  \item \textbf{Tissue}: cell type fingerprints distribute broadly across
    exclusion, layered, clustered, and hierarchical regimes, reflecting the
    diversity of tissue-scale organization.
\end{itemize}

44 of 72 empirical fingerprints (61\%) are flagged as out-of-regime
(NN distance exceeds the 95th percentile threshold of 0.260), suggesting
that real biological spatial organization often reflects mixtures or
mechanisms not captured by single-parameter simulations.

\begin{figure}[htbp]
  \centering
  \includegraphics[width=\textwidth]{../figures/fig_regime_affinity.pdf}
  \caption{KDE-based soft regime affinity profiles for all 72 empirical
  fingerprints. Each row shows the normalized probability of the fingerprint
  under each regime's distribution in PCA space.}
  \label{fig:regime_affinity}
\end{figure}

\subsection{Mixture recovery is limited}

Known mixtures of simulated regimes are poorly recovered by the KDE-based
affinity framework (Fig.~\ref{fig:mixture_validation}). A 70\% layered +
30\% clustered mixture yields max recovery error of 0.387; a 50\%
compartmental + 50\% exclusion mixture yields 0.999 (near-total failure).
This confirms that the soft affinity framework captures broad regime
proximity but cannot reliably decompose mixtures---consistent with the
null cloud degeneracy and the limitations of single-mechanism KDEs for
multi-modal patterns.

\begin{figure}[htbp]
  \centering
  \includegraphics[width=\textwidth]{../figures/fig_mixture_validation.pdf}
  \caption{Mixture validation: known mixture proportions (blue) versus
  recovered KDE-based affinities (coral) for three test mixtures.}
  \label{fig:mixture_validation}
\end{figure}


% ══════════════════════════════════════════════════════════════
\section{Discussion}
\label{sec:discussion}

\paragraph{What VMR fingerprints resolve.}
The central result is that VMR-derived descriptors define a real, stable
clustering manifold: mechanisms that produce above-Poisson variance
(clustering, compartmentalization, hierarchical nesting, branch
confinement) are distinguishable in a 6D descriptor space compressed
to 3 PCs explaining 96.0\% of variance. This is not a trivial result---it
means the VMR framework does more than detect ``something non-random'';
it provides mechanistic discrimination among clustering-positive processes.

\paragraph{What VMR fingerprints do not resolve.}
Equally important is the negative result: exclusion (sub-Poisson
regularity), layered organization, and randomness produce nearly
indistinguishable VMR fingerprints. This is because VMR primarily measures
count overdispersion, which is near 1 for all three. Discriminating these
regimes would require additional descriptors---anisotropy for layered
patterns, pair correlation functions for exclusion---that go beyond the
VMR curve.

\paragraph{Implications for empirical studies.}
The 61\% out-of-regime rate for empirical fingerprints is informative,
not alarming. It means that most real biological systems produce spatial
patterns more complex than any single simulated mechanism, likely
reflecting mixtures or multi-scale processes. The simulation-defined
regime space provides a coordinate system for locating empirical data,
not a classification scheme.

\paragraph{Limitations.}
Each regime is swept along a single parameter axis; real mechanisms have
multiple interacting parameters (e.g., cluster radius $\times$ number of
clusters $\times$ background fraction). Multi-parameter sweeps would
produce denser manifolds and potentially reveal additional structure.
The KDE-based mixture recovery fails because KDEs fitted to
single-mechanism clouds do not model the between-regime interpolation
that mixtures create. More sophisticated mixture models (e.g.,
density-ratio estimation) could improve this.

\paragraph{Relationship to companion studies.}
This work provides theoretical context for two companion empirical
studies: one applying the VMR framework to cell type organization in the
mouse brain (spatial transcriptomics), and another to protein
nanoclusters (SMLM). The regime space defined here offers a shared
reference frame for interpreting fingerprints across biological scales.


% ══════════════════════════════════════════════════════════════
\section*{Data Availability}

All code, simulated results, and figures are available at
\url{https://github.com/saivoru777-ship-it/biospatial-regimes}.


% ══════════════════════════════════════════════════════════════
\bibliographystyle{unsrtnat}
\bibliography{references}


% ══════════════════════════════════════════════════════════════
\clearpage
\appendix
\renewcommand{\thefigure}{S\arabic{figure}}
\setcounter{figure}{0}

\section*{Supplementary Figures}

\begin{figure}[htbp]
  \centering
  \includegraphics[width=\textwidth]{../figures/fig_pca_loadings.pdf}
  \caption{PCA loadings on the six fingerprint descriptors for each
  principal component.}
  \label{fig:pca_loadings}
\end{figure}

\begin{figure}[htbp]
  \centering
  \includegraphics[width=\textwidth]{../figures/fig_descriptor_distributions.pdf}
  \caption{Distributions of each fingerprint descriptor, colored by regime.}
  \label{fig:descriptor_distributions}
\end{figure}

\begin{figure}[htbp]
  \centering
  \includegraphics[width=\textwidth]{../figures/fig_separability.pdf}
  \caption{Regime separability: overall silhouette score (0.172) and
  pairwise Mahalanobis distances between regime clouds.}
  \label{fig:separability}
\end{figure}

\begin{figure}[htbp]
  \centering
  \includegraphics[width=\textwidth]{../figures/diag_pca_scatter.pdf}
  \caption{Diagnostic PCA scatter showing all 3,750 simulated descriptors
  colored by regime, with parameter sweep gradients visible as continuous
  clouds.}
  \label{fig:diag_pca}
\end{figure}

\begin{figure}[htbp]
  \centering
  \includegraphics[width=\textwidth]{../figures/diag_vmr_sweeps.pdf}
  \caption{VMR ratio curves across parameter sweeps for each regime,
  showing how fingerprint shape changes continuously with mechanism
  parameters.}
  \label{fig:diag_vmr}
\end{figure}

\begin{figure}[htbp]
  \centering
  \includegraphics[width=\textwidth]{../figures/diag_descriptor_boxplots.pdf}
  \caption{Box plots of each descriptor by regime, showing separation
  and overlap.}
  \label{fig:diag_boxplots}
\end{figure}

\begin{figure}[htbp]
  \centering
  \includegraphics[width=\textwidth]{../figures/diag_sweep_monotonicity.pdf}
  \caption{Sweep monotonicity: how each descriptor changes across the
  parameter sweep for each regime. Monotonic trends indicate interpretable
  parameter--descriptor relationships.}
  \label{fig:diag_monotonicity}
\end{figure}

\begin{figure}[htbp]
  \centering
  \includegraphics[width=\textwidth]{../figures/fig_pilot_branch_test.pdf}
  \caption{Pilot test confirming that branch-constrained patterns are
  separable from clustered patterns in the VMR descriptor space.}
  \label{fig:pilot}
\end{figure}

\end{document}
